La calibración sigue el enfoque de simulaciones de comportamiento aceptable \citep{beven1992}. Los criterios se fijaron antes de examinar los resultados:
\begin{enumerate}
\item fracción $R$ de puntos sobre calles con presencia informada por las fuentes en que $h_\mathrm{max}>h^*$, con $R\ge0{,}6$;
\item fracción $R_\mathrm{sit}$ de los cuatro sitios con presencia informada por las fuentes que el flujo alcanza dentro de \SI{40}{m}, con $R_\mathrm{sit}\ge3/4$;
\item llegada a menos de \SI{200}{m} del cruce de General Blanche con Padre Hurtado;
\item fracción $F_\mathrm{per}$ del área afectada que cae dentro del perímetro municipal ampliado en \SI{100}{m}, con $F_\mathrm{per}\ge0{,}5$. Desde el 11 de octubre, cuando se publicó la zona afectada oficial \citep{bbnn2026visor}, $F_\mathrm{per}$ se calcula con ese polígono y solo sobre el barro simulado en zona urbana (a menos de \SI{40}{m} de un edificio), para no contar los cerros mojados por la lluvia.
\end{enumerate}
El umbral es $h^*=\SI{5}{cm}$, con sensibilidad para \SI{2}{cm} y \SI{10}{cm}. Cada simulación aceptable recibe un peso informal $L=R\,R_\mathrm{sit}\,F_\mathrm{per}$. Las calles cerradas por el municipio se usan solo como contexto, porque un cierre puede ser preventivo. El índice de Jaccard respecto del perímetro se informa como dato descriptivo, ya que el perímetro es operativo y no una huella cartografiada.

El barrido recorre las hipótesis H1, H2 y H3; volúmenes de mezcla de \SI{80000}{\metre\cubed}, \SI{150000}{\metre\cubed} y \SI{300000}{\metre\cubed}; caudales pico de \SI{100}{\metre\cubed\per\second} y \SI{250}{\metre\cubed\per\second}; y tres muestras reológicas (Glenwood~1, Glenwood~2b y Aspen~2). El cribado usó CFL~1. Se completaron controles con paso temporal fino (z01, z02, z26, z29 y z32) y el ensayo con malla fina en el campus (m36), según las secciones~\ref{sec:res-verif} y~\ref{sec:res-posterior}.
